Trading agents built on large language models (LLMs) act on retrieved news, so an attacker who edits or injects news can change their trades without touching prices or code. We present \rhal{}, a defense layer for such agents and a framework to evaluate it. The layer grounds the agent's reading of the news in market evidence that the attacker cannot edit: trailing price and microstructure features and the realized outcomes of similar past windows, all computed without lookahead. A risk layer then gates and sizes each decision. We test it against an adaptive misinformation attacker in the style of AutoRedTrader and a frozen replay of its attack, over \val{NseedsRun} paired simulated markets. Grounding lowers misinformation-driven attack success (net ASR) from \val{CmpNetNone}\% to \val{CmpNetGrounded}\%, at a cost of \val{CmpCostGrounded}\pp{} of clean-market return. It also outperforms detection filtering and robust aggregation of the news. The gain persists when a local 7B LLM reads the news, on \val{NseedsReal} windows of real BTC-USD prices, against an attacker that aligns its edits with the market evidence, and when the misinformation also moves prices. A hold-on-conflict risk rule backfires because the attacker controls the conflict it reacts to; shrinking the trade instead removes the penalty. An attacker that maximizes text--market conflict is the strongest we found: it raises net ASR against the full layer by \val{CmpAwareAdaptiveConflictVsStdFull}\pp{}. Placebo controls credit the defense only for misinformation-driven flips: \val{ShareNoiseOfNaive}\% of naive attack success also appears when harmless paraphrases are injected.
